\section{Introduction}
\label{sec:intro}

% --- Background (reuse from the previous version of the paper) ---
Physics-informed neural networks (PINNs)~\cite{raissi2019} solve forward and inverse
partial differential equation (PDE) problems by embedding the governing equation, boundary
conditions and, for inverse problems, observations into the neural network's training loss.
We study a one-dimensional temperature diffusion problem, the seawater temperature equation of
\citet{han2026}, under Dirichlet, Neumann and Robin boundary conditions.
Fast, differentiable surrogates for this equation are attractive for energy-aware control of
engineered water bodies, such as swimming pools or other managed basins, where a model that
returns temperature and its gradients cheaply can drive heating, circulation or shading decisions
in an optimal-control or model-predictive-control loop.

Hybrid quantum--classical PINNs (QCPINNs), in which part of the network is replaced by a variational
quantum circuit (VQC), have been reported to match or exceed classical PINNs while using far fewer
trainable parameters~\cite{farea2025}. Parameter count has now become the standard efficiency metric
in this literature.

% --- Gap ---
\paragraph{Gap.} Wall-clock cost is rarely the headline metric and, where it is reported, it is a
single aggregate figure (see Related work below). Yet near-term VQCs are carried out using classical
simulation, since its cost increases with circuit width and depth (Section~\ref{sec:background}).
Training a PINN also requires first- and second-order derivatives of the network with respect to its
\emph{inputs}, which are more expensive than a plain forward pass and interact with the simulated
circuit.

\paragraph{Research questions.}
\begin{description}
  \item[RQ1] How large is the parameter-count advantage of hybrid models, and where does it come from?
  \item[RQ2] What timing cost accompanies it, and how does that cost scale with qubit count, circuit
  depth, batch size and simulator backend?
\end{description}

% --- Contributions ---
\paragraph{Contributions.}
\begin{enumerate}
  \item A parameter-count breakdown of a classical PINN and four hybrid configurations under a common
  accounting.
  \item Stage-wise timing of the forward pass, the first-order input gradient and the second-order
  input gradient.
  \item A sweep of qubit count, circuit depth, batch size and simulator backend.
\end{enumerate}
Accuracy comparisons, gradient-norm ablations and barren-plateau discussion are not the subject of
this paper.

\paragraph{Outline.} Section~\ref{sec:background} defines the PINN loss, the timed stages and the VQC basics
used throughout; Section~\ref{sec:models} describes the classical baseline and the four hybrid configurations;
Section~\ref{sec:setup} details the timing protocol and the size-matched classical controls. This draft covers
the setup only -- the parameter-count and timing results, discussion and conclusion follow in later sections
of the full paper.

\paragraph{Related work.}
Hybrid quantum PINNs have been proposed for fluid flow, wave, Helmholtz and Klein--Gordon
problems~\cite{farea2025,leong2025}, reservoir seepage~\cite{rao2025}, hydrological
forecasting~\cite{hewage2026} and Maxwell's equations~\cite{chen2025}, with theoretical work on universal
approximation and barren-plateau mitigation~\cite{lantigua2025}. The headline efficiency claim in this literature is
parameter count: QCPINNs are reported to reach comparable accuracy with 10--30\% of the parameters of a classical
PINN~\cite{farea2025}, hybrid models with fewer parameters than the best classical model~\cite{leong2025,trahan2024},
and about 44\% fewer parameters together with roughly three times fewer training epochs~\cite{hewage2026}.

Cost is acknowledged but usually reported as one aggregate figure. \citet{farea2025} report an increase of 15- to
40-fold in time per training iteration relative to classical PINNs, at similar memory. \citet{leong2025} state that
simulating the quantum layers imposes an overhead of at least two orders of magnitude for the same number of
trainable parameters, but compare training efficiency through epochs rather than wall-clock time. \citet{trahan2024}
find a quantum PINN to take nearly four times as long as a classical PINN with a similar number of parameters on a
2D Poisson problem.
These figures differ widely and depend on hardware, simulator, circuit and batch size; none of them, to our
knowledge, breaks the cost down into forward, first-order and second-order input-gradient stages, or varies qubit
count, depth, batch size and backend in a controlled way. On the simulation side, \citet{chen2025} develop a
GPU-accelerated simulation library for computing circuit outputs and their derivatives, to make QPINN training
practical. This paper fills the gap with a controlled timing study alongside the parameter count.
